\section{ANEXO G. DESPLIEGUE DE LA INFRAESTRUCTURA EN LA NUBE
(RAILWAY)}\label{anexo-g.-despliegue-de-la-infraestructura-en-la-nube-railway}

Para validar la operatividad de EVMAudit en un entorno accesible y
simular un escenario de producción real, se ha procedido al despliegue
de la arquitectura contenedorizada en la plataforma de Plataforma como
Servicio (PaaS) \textbf{Railway}. A continuación, se detallan las
especificaciones del entorno, las restricciones técnicas de hardware
identificadas y las optimizaciones de ingeniería aplicadas en el código
fuente para garantizar la estabilidad del sistema.

\subsection{ANEXO G.1. Aprovisionamiento y Configuración del Entorno
Cloud}\label{anexo-g.1.-aprovisionamiento-y-configuraciuxf3n-del-entorno-cloud}

El proceso de despliegue en la infraestructura de la nube se ha
estructurado bajo las siguientes directrices operativas:

\begin{itemize}
\item
  \textbf{Selección del Nivel de Servicio:} La instancia se ha
  instanciado haciendo uso del nivel gratuito (\emph{Free Tier}) de la
  plataforma, el cual provee un crédito base de \$5 USD o un límite
  temporal de 30 días de cómputo.
\item
  \textbf{Aislamiento Regional:} Con el objetivo de minimizar la
  latencia de red en las peticiones HTTP y optimizar la transferencia de
  datos, se ha seleccionado la región europea con nodo central en
  \textbf{Ámsterdam (EU)}.
\item
  \textbf{Mapeo de Infraestructura y Orquestación:} El aprovisionamiento
  se realiza directamente vinculando el contenedor web a la imagen
  Docker compilada y almacenada en el registro de GitHub (ghcr.io),
  exponiendo de manera transparente la API del \emph{backend}
  desarrollada en FastAPI.
\item
  \textbf{Enrutamiento y Capa de Enlace (SSL):} La plataforma genera de
  manera dinámica un nombre de dominio completamente cualificado (FQDN)
  provisto de seguridad criptográfica TLS/SSL (HTTPS) para el acceso
  público a la interfaz de usuario:
  \url{https://evmaudit-production.up.railway.app/}.
\end{itemize}

\subsection{ANEXO G.2. Limitaciones de Hardware y el Problema del
Desbordamiento de Memoria
(OOM)}\label{anexo-g.2.-limitaciones-de-hardware-y-el-problema-del-desbordamiento-de-memoria-oom}

La modalidad gratuita de la plataforma PaaS impone restricciones
estrictas sobre los recursos de hardware asignados a cada contenedor,
parametrizados de la siguiente forma:

\begin{itemize}
\tightlist
\item
  \textbf{Capacidad de Cómputo (CPU):} 2 vCPU virtuales compartidas.
\item
  \textbf{Memoria Volátil (RAM):} 1 GB con un límite estricto de cuota
  (\emph{Plan Limit}).
\end{itemize}

Bajo un escenario de despliegue convencional en servidores dedicados o
infraestructura local, el pipeline de análisis de EVMAudit se ejecuta
sin restricciones debido a la disponibilidad de memoria elástica. Sin
embargo, en el entorno de la nube restringido, el motor de
\emph{fuzzing} basado en propiedades \textbf{Echidna} presenta un
problema crítico de arquitectura.

Echidna, al estar desarrollado en Haskell, requiere de forma nativa una
reserva inicial y un espacio de intercambio que supera con creces el
gigabyte de memoria RAM para gestionar el mapa de cobertura del binario
y la generación de casos de prueba. Al alcanzar el umbral crítico de 1
GB asignado por Railway, el demonio del sistema operativo anfitrión
(\emph{Kernel Out-of-Memory Killer}) destruía de manera abrupta el
contenedor para salvaguardar la integridad del nodo, provocando la caída
del servicio web y reportando un error de tipo \emph{OOM}.

\subsection{ANEXO G.3. Optimización del Sistema en Tiempo de Ejecución
(RTS) de
Haskell}\label{anexo-g.3.-optimizaciuxf3n-del-sistema-en-tiempo-de-ejecuciuxf3n-rts-de-haskell}

Para mitigar el desbordamiento de memoria sin alterar las capacidades
analíticas esenciales de la herramienta, se realizó una intervención a
nivel de código en el módulo de control del \emph{pipeline}
(run\_echidna). La solución consistió en inyectar directivas específicas
orientadas a reconfigurar los parámetros del \textbf{Runtime System
(RTS)} del compilador de Glasgow Haskell (GHC) empaquetados dentro del
binario de Echidna.

La estructura de invocación del proceso fue modificada e implementada en
Python mediante el siguiente diseño de argumentos:

\begin{Shaded}
\begin{Highlighting}[]
\NormalTok{command }\OperatorTok{=}\NormalTok{ [}
    \StringTok{"echidna"}\NormalTok{,}
\NormalTok{    contract\_path,}
    \StringTok{"{-}{-}contract"}\NormalTok{, contract\_name,}
    \StringTok{"{-}{-}format"}\NormalTok{, }\StringTok{"json"}\NormalTok{,}
    
    \CommentTok{\# Reducción del espacio de búsqueda del Fuzzer}
    \StringTok{"{-}{-}test{-}limit"}\NormalTok{, }\StringTok{"100"}\NormalTok{,  }
    
    \CommentTok{\# Optimización de la compilación inicial}
    \StringTok{"{-}{-}solc{-}args"}\NormalTok{, }\StringTok{"{-}{-}optimize{-}runs 0"}\NormalTok{, }
    
    \CommentTok{\# Apertura de las opciones del Runtime System de Haskell}
    \StringTok{"+RTS"}\NormalTok{,      }
    
    \CommentTok{\# Parámetros estrictos de control de memoria y concurrencia}
    \StringTok{"{-}M950m"}\NormalTok{,    }
    \StringTok{"{-}c"}\NormalTok{,        }
    \StringTok{"{-}N1"}\NormalTok{,       }
    
    \CommentTok{\# Cierre de las opciones RTS}
    \StringTok{"{-}RTS"}       
\NormalTok{]}
\end{Highlighting}
\end{Shaded}

A continuación se expone la justificación técnica detrás de los
modificadores inyectados:

\begin{itemize}
\tightlist
\item
  \textbf{Restricción de Ensayos (-\/-test-limit 100):} Al parametrizar
  el límite de pruebas en 100 iteraciones, se acota la profundidad del
  grafo de ejecución generado por el fuzzer. Esto reduce el consumo de
  memoria acumulativo a lo largo del tiempo de computación.
\item
  \textbf{Techo Infranqueable de Memoria (-M950m):} Esta directiva
  establece que el asignador de memoria de Haskell tiene prohibido
  estrictamente reclamar más de 950 megabytes del espacio de usuario. Al
  situar este límite ligeramente por debajo del gigabyte real de
  Railway, se evita que el sistema operativo de la nube elimine el
  proceso por exceder la cuota física.
\item
  \textbf{Recolección de Basura Agresiva (-c):} Activa un algoritmo de
  recolección de residuos más severo en el recolector de basura de
  Haskell (\emph{Garbage Collector}). En lugar de acumular objetos
  intermedios en la memoria RAM, el sistema libera los recursos
  obsoletos inmediatamente después de cada evaluación de propiedad.
\item
  \textbf{Concurrencia Confinada (-N1):} Limita la ejecución del entorno
  de ejecución a un único hilo de procesamiento, evitando la duplicación
  de estructuras de datos en memoria asociadas al paralelismo de hilos
  nativos.
\end{itemize}

Es importante destacar que la optimización descrita no se integró d
forma estática en la construcción del archivo Dockerfile (lo que habrí
alterado el comportamiento de la imagen base de manera permanente), sin
que se aplicó directamente sobre el código fuente desplegado en el
contenedor en ejecución (\emph{runtime}). Bajo condiciones de despliegue
convencionales en infraestructuras con escalabilidad elástica o recurso
dedicados de hardware, esta intervención técnica resultaría
completamente innecesaria, ya que la aplicación contaría con la memoria
suficiente para procesar el \emph{pipeline} por defecto. Por
consiguiente, esta modificación responde de manera estricta a un
mecanismo de mitigación ad hoc, implementado exclusivamente para sortear
las limitaciones físicas del entorno gratuito y evitar la interrupción
forzada del servicio web por falta de memoria.

\textbf{Conclusión del Despliegue:} La implementación de estas
salvaguardas de bajo nivel ha permitido que la aplicación web de
\textbf{EVMAudit} opere de manera completamente estable y fluida en la
nube. Pese a las severas restricciones del entorno gratuito de Railway,
el sistema es capaz de completar con éxito el pipeline completo de
auditoría en siete pasos sin registrar caídas en el servicio ni
excepciones por falta de recursos.
